% Figure: Doppler frequency ratio versus source Mach number for an
% approaching and a receding source, stationary observer.
% f'/f0 = 1 / (1 -+ M) for source moving toward (-) / away (+) the observer.
\begin{figure}[htbp]
\centering
\begin{tikzpicture}
\begin{axis}[
  width=12cm, height=7cm,
  xlabel={source speed as fraction of sound speed, $M=v_s/c$},
  ylabel={received-to-emitted frequency ratio $f'/f_0$},
  xmin=0, xmax=0.9, ymin=0, ymax=6,
  axis lines=left,
  domain=0:0.9, samples=200,
  legend pos=north west,
  legend cell align=left,
  legend style={font=\scriptsize},
  tick label style={font=\scriptsize},
  label style={font=\small},
  grid=major, grid style={enggray!20},
]
% approaching: f'/f0 = 1/(1 - M)
\addplot[engblue, very thick] {1/(1-x)};
\addlegendentry{source approaching}
% receding: f'/f0 = 1/(1 + M)
\addplot[engred, very thick] {1/(1+x)};
\addlegendentry{source receding}
% unison reference
\addplot[enggray, thin, dashed] {1};
\node[enggray, font=\scriptsize, anchor=south west] at (axis cs:0.02,1.05) {$f'=f_0$};
\end{axis}
\end{tikzpicture}
\caption[Doppler frequency ratio versus source Mach number]{Received
frequency relative to emitted frequency for a stationary observer and a
source moving at speed $v_s$, written against the source Mach number
$M=v_s/c$. An approaching source raises the pitch as $1/(1-M)$, which
diverges as the source nears the sound speed; a receding source lowers
it as $1/(1+M)$, bounded below by one-half at $M\to 1$. The asymmetry
between the two branches is the audible content of a passing siren: the
upward shift on approach is larger than the downward shift on recession,
so the drop heard as the source passes is the sum of the two and is
never symmetric about the rest pitch. The curves assume motion directly
along the line of sight; off-axis motion replaces $v_s$ by its
line-of-sight component, which sweeps through zero at the closest
approach and produces the characteristic glide of a passing vehicle.}
\label{fig:vol03ch07:doppler}
\end{figure}
